\chapter{Cardiopulmonary Resuscitation}
\index{Cardiopulmonary Resuscitation}

\section*{Definition and Overview}
Cardiopulmonary resuscitation (CPR) is the emergency technique used to preserve blood flow to the brain and other vital organs when a person's heart has stopped pumping effectively, a situation called cardiac arrest. It combines rhythmic chest compressions that manually pump the heart, with rescue breaths that provide oxygen to the lungs. CPR does not restart a stopped heart on its own, but it buys time by keeping oxygen-rich blood flowing until a defibrillator can deliver a shock to restore a normal heart rhythm, or until emergency services arrive. It is one of the few truly life-saving interventions a bystander can perform, and it is safe to attempt in most situations.

\section*{Epidemiology}
Cardiac arrest is a leading cause of sudden death, affecting roughly one in 2000 adults each year outside hospital in high-income countries. In many countries this means tens of thousands of events each year, most occurring in the home. Survival from out-of-hospital cardiac arrest is generally low, often under 10\%, but it increases several-fold when a bystander starts CPR immediately and a defibrillator is used early. For every minute without CPR and defibrillation, the chance of survival falls by about 10\%, which makes rapid, competent bystander action the single most important factor in survival after public recognition and the emergency call.

\section*{Aetiology and Pathophysiology}
The conditions that lead to cardiac arrest, and therefore to the need for CPR, are important to understand:
\begin{itemize}
    \item \textbf{Ventricular fibrillation:} the commonest initial rhythm in adults, in which the heart's lower chambers quiver chaotically and pump no blood; this rhythm responds to defibrillation
    \item \textbf{Coronary artery disease:} heart attacks and severe blockages that destabilise the heart's electrical system
    \item \textbf{Cardiomyopathy, valve disease, and heart failure:} structural problems that predispose to dangerous rhythms
    \item \textbf{Inherited arrhythmia syndromes:} conditions such as long QT syndrome and Brugada syndrome, which may affect young, otherwise healthy people
    \item \textbf{Other causes:} drowning, drug overdose, severe blood loss, low blood oxygen, electrolyte disturbances, and severe hypothermia
\end{itemize}
When the heart stops pumping, blood flow to the brain ceases and consciousness is lost within seconds. Chest compressions create artificial circulation, maintaining blood flow to the heart and brain until the underlying rhythm disturbance can be treated.

\section*{Clinical Features}
The need for CPR is recognised by the sudden, unmistakable signs of cardiac arrest:
\begin{itemize}
    \item Sudden collapse and loss of responsiveness
    \item Not breathing, or only gasping (agonal breathing), which is not effective breathing
    \item No pulse and no signs of life
    \item Pale, grey, or blue skin and lips
    \item Some people have warning symptoms in the hour beforehand, such as chest pain, palpitations, dizziness, or breathlessness
\end{itemize}

\section*{Differential Diagnosis}
Before starting CPR, a collapsed and unresponsive person must be quickly assessed, because several conditions mimic cardiac arrest:
\begin{itemize}
    \item Simple fainting (syncope), in which breathing and a pulse continue and recovery is usually rapid
    \item Seizures, which often recover with breathing present, but can occur during cardiac arrest itself
    \item Stroke causing sudden collapse and unconsciousness
    \item Severe hypoglycaemia in people with diabetes
    \item Drug or alcohol overdose causing profound unconsciousness
    \item Choking, where a blocked airway stops breathing
\end{itemize}
In any unresponsive person who is not breathing normally, assume cardiac arrest and start CPR immediately; waiting to confirm the cause costs vital time.

\section*{When to Seek Medical Care}
Cardiac arrest is always a 000 emergency, and CPR is performed while emergency help is on the way. Do not delay the emergency call to assess further.
\textbf{Contact your local emergency services for:}
\begin{itemize}
    \item A person who has collapsed and is unresponsive
    \item Not breathing or only gasping
    \item Sudden severe chest pain, collapse, or loss of consciousness
    \item While waiting for help, start CPR and ask someone to fetch an automated external defibrillator (AED) if one is nearby
\end{itemize}

\section*{Diagnosis and Investigations}
Cardiac arrest is diagnosed clinically, at the bedside, by unresponsiveness and abnormal or absent breathing. Diagnosis and the investigations that follow focus on finding the cause and guiding treatment:
\begin{itemize}
    \item \textbf{Continuous cardiac monitoring:} during resuscitation, to identify shockable rhythms such as ventricular fibrillation
    \item \textbf{Capnography:} measures exhaled carbon dioxide to confirm that compressions and ventilation are effective
    \item \textbf{Blood tests:} electrolytes, troponin, glucose, and blood gases to identify reversible causes
    \item \textbf{Electrocardiogram (ECG):} after stabilisation, to detect heart attack patterns or inherited rhythm syndromes
    \item \textbf{Echocardiogram:} to assess heart function and structural damage
    \item \textbf{Imaging:} CT of the head or chest where a stroke, pulmonary embolism, or aortic problem is suspected
    \item \textbf{In survivors without an obvious cause:} assessment for inherited cardiac conditions, including screening of close family members
\end{itemize}

\section*{Management}
The management of cardiac arrest follows a standard, time-critical sequence known as the chain of survival:
\begin{itemize}
    \item \textbf{Call for help:} dial 000 immediately, put the phone on speaker, and follow the dispatcher's instructions
    \item \textbf{CPR:} push hard and fast in the centre of the chest, at about 100--120 compressions per minute, allowing the chest to fully recoil between compressions; give rescue breaths if trained and willing
    \item \textbf{Defibrillation:} apply an AED as soon as it arrives and follow its voice prompts; deliver a shock if advised, then resume CPR immediately
    \item \textbf{Advanced life support:} paramedics and hospital teams use airway devices, drugs to support the heart and rhythm, and treat reversible causes
    \item \textbf{Care after return of circulation:} targeted temperature management, treatment of the underlying cause, and care in an intensive care unit
    \item \textbf{Rehabilitation:} survivors may need physical, cognitive, and psychological support, and an implantable defibrillator where arrhythmia risk remains high
\end{itemize}

\section*{Complications}
Even when CPR is performed correctly, it can cause injuries, and cardiac arrest itself carries serious after-effects:
\begin{itemize}
    \item Fractures of the ribs or breastbone from chest compressions
    \item Brain injury from oxygen deprivation, ranging from subtle memory problems to permanent disability
    \item Lung damage, including aspiration of stomach contents into the lungs
    \item Injury to the liver or other organs, which is rare
    \item Kidney, liver, or other organ failure from low blood flow during the arrest
    \item Post-cardiac-arrest psychological distress, including anxiety, depression, and post-traumatic stress disorder
    \item Recurrence of cardiac arrest if the underlying cause is not corrected
\end{itemize}
These risks are far outweighed by the benefit of CPR in a person in cardiac arrest, who will otherwise die.

\section*{Prognosis and Outcomes}
The outlook depends strongly on how quickly treatment begins. Survival is far more likely when CPR is started within minutes and a shock is delivered promptly, and for every minute without CPR and defibrillation the chance of survival falls substantially. In many out-of-hospital arrests the outcome remains poor, but bystander CPR and public-access defibrillators improve survival several-fold. Among survivors, recovery of brain function varies widely: some people return to a full and active life, while others have lasting disability. In-hospital arrests and those with an initial shockable rhythm generally carry a better outlook.

\section*{Prevention and Self-Care}
\begin{itemize}
    \item Learn CPR and how to use an AED, and refresh your skills every one to two years
    \item Recognise the early warning signs of a heart attack, such as chest discomfort, and seek help immediately
    \item Treat high blood pressure, high cholesterol, and diabetes under medical supervision
    \item Do not smoke, and limit alcohol intake
    \item Follow specialist advice for inherited heart conditions, including screening of family members
    \item Keep AEDs accessible and maintained at workplaces, sports grounds, and community venues
\end{itemize}

\section*{Sources and Further Reading}
The following sources provide authoritative, up-to-date information on cardiopulmonary resuscitation and training.
\begin{itemize}
    \item St John Ambulance Australia. \textit{CPR and first aid training.} https://www.stjohn.org.au/
    \item Australian Resuscitation Council. \textit{CPR guidelines and chain of survival.} https://www.resus.org.au/
    \item NHS. \textit{First aid: how to perform CPR.} https://www.nhs.uk/conditions/first-aid/
    \item American Heart Association. \textit{CPR and emergency cardiovascular care.} https://www.heart.org/
    \item MedlinePlus. \textit{Cardiopulmonary resuscitation (CPR).} https://medlineplus.gov/cpr.html
\end{itemize}
